\section{\label{sec:preparation}Data Preparation}

testing

\subsection{\label{sec:conversion}Conversion of a Digraph}

\justifying
Once our Terraform infrastructure has been output as a directed graph, we need to bring it in to Python so we can
perform operations on it. To run Terraform commands in an automated fashion,
the \href{https://pypi.org/project/python-terraform/}{python-terraform module} is used.
\vspace{2mm}

\justifying
Directed graphs are ingested into Python as an object using the \href{https://pypi.org/project/networkx/}{networkx module}. Now we are prepared
to perform computational operations on the digraph.
\vspace{2mm}

\justifying
\begin{figure}[H]
	\includegraphics[width=12cm]{graph3.png}
	\caption{Directed graph as ingested by Networx}
	\label{digraph}
\end{figure}

\justifying
Note that new numerical labels are displayed in place of the original labels. These are
stored in the node objects in case we need them later. The root node of the adjacency matrix is labeled
as the 0th element in the graph. Also of note is that the graphs displayed in figure \ref{pygraph} and figure \ref{digraph} can be described as \href{https://en.wikipedia.org/wiki/Graph\_isomorphism}{isomorphic}.
\vspace{2mm}

\justifying
Cleaning data using the \href{}{pyjanitor module}.
\vspace{2mm}

\subsection{\label{sec:adjacency}The Adjacency Matrix}

\justifying
The adjacency matrix is a mathematical representation of our graph. The nodes are represented in a
columnar format known as a matrix. Continuing with our previous example, the representation of our

Add the identity matrix to the adjacency matrix using an update rule.
The goal is to prevent the central node from being excluded. Doing this without increasing the scale requires something like multiplication by the inverse of the degree matrix, aka ``mean pooling''. The
most popular graph convolutional layer is the GCN update rule. We will not be able to attach complex features to the edges.

\subsection{\label{sec:embedding}Node Embedding}

\justifying
Assigning weights and features to the edges of the graph. Labeling.

Attach feature vectors to nodes.

Compute arbitrary message vectors between neighbor nodes.

\subsection{\label{sec:conv}Convolutional Layers}